% ------------------------------------------------------------------------------
%  Capitolo 1 – Business Understanding
% ------------------------------------------------------------------------------
\chapter{Business Understanding}
\label{ch:business_understanding}

% -- 1.1 ----------------------------------------------------------------------
\section{Obiettivi di Business}
\label{sec:obiettivi_business}

A fronte dei limiti operativi emersi nella fase di assistenza ai pazienti in attesa e delle nuove funzionalità di tracciamento introdotte nella piattaforma, il valore fondamentale che l'ecosistema SINTONIA deve generare per l'ASL Campania consiste nel trasformare il periodo di attesa clinica in una risorsa attiva. L'obiettivo primario è capitalizzare l'interazione quotidiana dell'utente con l'applicativo per strutturare un sistema di monitoraggio continuo, capace di supportare concretamente l'azione preventiva e l'efficacia del Triage psicologico.

Declinando questo valore sui singoli moduli previsionali, si definiscono i seguenti obiettivi aziendali:

\begin{itemize}
    \item \textbf{Gestione delle Emergenze:} Il sistema deve permettere agli psicologi incaricati di ricevere segnalazioni immediate in caso di espressione di intenti critici o autolesionistici così da garantire un intervento di emergenza tempestivo e ridurre drasticamente il rischio di eventi avversi durante l'attesa.

    \item \textbf{Verifica Clinica e Allocazione Risorse:} Il sistema deve permettere all'ASL Campania di confrontare costantemente l'andamento emotivo reale del paziente con i punteggi dichiarati nei questionari clinici periodici così da identificare tempestivamente discrepanze, focalizzare l'attenzione medica sui casi a maggior rischio silente e ottimizzare di conseguenza le priorità di presa in carico.

    \item \textbf{Mantenimento del Canale Assistenziale:} Il sistema deve permettere alla piattaforma stessa di individuare e contrastare proattivamente i cali di interesse e interazione da parte degli utenti così da preservare un livello di coinvolgimento costante, garantendo il flusso continuo di informazioni necessario al corretto funzionamento dell'intero apparato di Triage.
\end{itemize}

È fondamentale sottolineare le aspettative e i vincoli operativi condivisi con gli \textit{stakeholder} sanitari in merito all'adozione di questi strumenti. In primo luogo, il sistema assume esclusivamente un ruolo di supporto decisionale: le indicazioni fornite non costituiscono in alcun caso una diagnosi automatizzata e non si sostituiscono al giudizio clinico e alla responsabilità del professionista. In secondo luogo, il volume delle notifiche e degli avvisi generati deve mantenersi rigorosamente entro limiti operativamente sostenibili, per evitare di sovraccaricare il personale già soggetto a vincoli di organico. Infine, gli esiti prodotti dai moduli previsionali devono essere presentati attraverso interfacce estremamente intuitive e di immediata interpretazione, affinché possano essere impiegati e compresi dal personale sanitario senza richiedere alcuna specifica competenza tecnica o informatica.

% -- 1.2 ----------------------------------------------------------------------
\section{Criteri di Successo e Metriche di Valutazione}
\label{sec:criteri_successo}

Dopo aver delineato gli obiettivi del progetto e le funzionalità previste (Capitolo~\ref{ch:introduzione}), è necessario stabilire criteri concreti per valutare se tali obiettivi saranno effettivamente raggiunti. Questa sezione definisce due livelli complementari di valutazione: i \textbf{Criteri di Successo} (\textit{Business Success Criteria}), ossia i KPI clinico-operativi con cui l'ASL Campania misurerà l'impatto dei componenti sul flusso di lavoro e sulla qualità dell'assistenza, e i corrispondenti \textbf{Obiettivi Tecnici}, ossia le soglie metriche che guideranno sviluppo e validazione dei componenti. La separazione tra i due piani assicura che ogni miglioramento nelle performance tecniche si traduca in un beneficio tangibile per pazienti e operatori sanitari.

\subsection{Criteri di Successo}
\label{subsec:business_success}

I criteri di successo rappresentano le metriche operative e cliniche attraverso cui l'ASL Campania potrà valutare l'efficacia dell'integrazione dei componenti di Intelligenza Artificiale nella piattaforma SINTONIA. Per ciascun componente descritto nella Sezione~\ref{sec:funzionalita_implementate}, vengono identificati i KPI più significativi.

\subsubsection{Sistema di Alert Suicidario}
\label{subsubsec:kpi_alert}

Per il componente di \textit{Red Flag Detection} (Sezione~\ref{subsec:red_flag}), i criteri di successo misurano il valore operativo che l'aggiunta dell'analisi del diario come canale di sorveglianza produce.

\begin{itemize}
    \item \textbf{Estensione della finestra di sorveglianza:} la copertura del monitoraggio del rischio suicidario, intesa come rapporto tra i giorni in cui il paziente è coperto da almeno un canale di sorveglianza e il totale dei giorni nel periodo di attesa, deve passare dall'attuale \textbf{7\%} (1 giorno ogni 14, corrispondente ai soli giorni di somministrazione del questionario) ad almeno il \textbf{50\%}. L'analisi del diario rende possibile questo salto poiché una pagina può essere scritta in qualsiasi giorno, a prescindere dal calendario dei questionari, estendendo la rete di sicurezza da un evento periodico a un monitoraggio distribuito lungo l'intero periodo di attesa.

    \item \textbf{Riduzione dei tentativi di suicidio durante il periodo di attesa:} la letteratura evidenzia che fino a \textbf{1 paziente su 6} tenta il suicidio mentre è in lista d'attesa per un trattamento psicologico \cite{we_need_to_talk2014}, dato che riflette l'assenza di strumenti di sorveglianza attiva nei periodi non coperti da contatto clinico. Poiché il 75\% delle persone decedute per suicidio ha negato intenti suicidari nell'ultimo contatto con un professionista sanitario \cite{busch2003clinical}, il diario personale --- in cui il paziente si esprime liberamente e senza interlocutore --- rappresenta un canale potenzialmente più rivelatore rispetto ai questionari a risposta chiusa. L'obiettivo è ridurre di almeno il \textbf{30\%} gli episodi critici (tentativi di suicidio, accessi al pronto soccorso per crisi psichiatriche, segnalazioni d'emergenza da parte di familiari) tra i pazienti in attesa, intercettando tempestivamente i segnali di rischio attraverso l'analisi del diario e attivando l'intervento clinico prima che la crisi si manifesti.
\end{itemize}

\subsubsection{Depression Prediction}
\label{subsubsec:kpi_phq9}

I criteri di successo per il componente di \textit{Depression Prediction} (Sezione~\ref{subsec:depression}) misurano il valore operativo prodotto dal sistema nel migliorare il processo di verifica della validità clinica dei dati raccolti.

\begin{itemize}
    \item \textbf{Copertura della verifica di validità:} la percentuale di compilazioni PHQ-9 sottoposte a controlli sistematici di congruenza deve passare dall'attuale copertura parziale e non deterministica (revisione a campione casuale, stimabile intorno al \textbf{10--15\%} del totale) al \textbf{100\% dei questionari valutabili}. Per \textit{valutabilità} si intende la preesistenza di una densità minima di input comportamentali (stati d'animo registrati) sufficienti ad alimentare una predizione affidabile. Su questi casi, il sistema garantisce che nessuna compilazione sospetta possa sfuggire all'attenzione degli psicologi, eliminando la dipendenza dalla casualità dell'ispezione manuale.

    \item \textbf{Riduzione delle compilazioni inattendibili nel flusso clinico:} la letteratura sul \textit{careless responding} stima che il \textbf{10--15\%} dei rispondenti a questionari autosomministrati produce risposte disattente, casuali o deliberatamente falsate \cite{meade2012careless}, percentuale sufficiente a introdurre bias significativi nelle decisioni cliniche. Nel contesto di SINTONIA, compilazioni PHQ-9 incongruenti possono alterare il punteggio di priorità del triage dinamico, facendo salire o scendere indebitamente un paziente nella lista d'attesa. L'obiettivo è dimezzare la quota di compilazioni inattendibili che entrano nel flusso decisionale clinico \textit{senza essere segnalate allo psicologo}, portandola dal \textbf{10--15\%} stimato dalla letteratura a meno del \textbf{5\%}: il sistema evidenzia automaticamente le discrepanze significative tra punteggio predetto e punteggio dichiarato, sottopone i casi sospetti all'attenzione dello psicologo, il quale valuta e decide in ultima istanza se confermare o escludere l'incongruenza.
\end{itemize}

\subsubsection{Prevenzione dell'Abbandono (Churn)}
\label{subsubsec:kpi_churn}

Per il componente di \textit{Churn Prevention} (Sezione~\ref{subsec:churn}), il criterio di successo primario è la \textbf{User Retention}: l'obiettivo è mantenere l'utilizzo attivo dell'applicazione da parte di almeno il \textbf{75\%} dei pazienti durante l'intero periodo di attesa precedente la presa in carico, dove per \textit{utilizzo attivo} si intende l'inserimento di almeno un'interazione significativa (diario, stato d'animo o questionario) nell'arco di ogni finestra di 7 giorni.

La letteratura scientifica evidenzia che le applicazioni mobile per la salute mentale presentano tassi di abbandono estremamente elevati: studi recenti riportano che fino al 50--80\% degli utenti cessa l'utilizzo entro le prime due settimane dall'installazione \cite{baumel2019mhealth}. In questo contesto, un obiettivo di retention del 75\% sull'intero periodo di attesa rappresenta un target ambizioso ma clinicamente necessario, e costituisce il benchmark rispetto al quale misurare l'efficacia delle strategie di \textit{nudging} prodotte dall'algoritmo genetico. La continuità dei dati comportamentali raccolti è inoltre un prerequisito per l'efficacia del triage dinamico e degli altri componenti predittivi: un tasso di abbandono elevato ridurrebbe progressivamente la quantità e la rappresentatività dei dati disponibili, degradando la qualità delle predizioni e della prioritizzazione clinica.

\subsection{Obiettivi Tecnici: Metriche di Valutazione}
\label{subsec:metriche_ml}

Gli obiettivi tecnici traducono i criteri di successo in metriche matematiche formali, ciascuna calibrata sulla tipologia di task affrontato dal rispettivo componente e sul livello di rischio clinico associato agli errori di predizione. Le soglie definite di seguito dovranno essere raggiunte su un test set indipendente, estratto mediante campionamento stratificato al fine di preservare la distribuzione originale delle classi e garantire la validità statistica della valutazione.

\subsubsection{Rischio Suicidario --- Classificazione Binaria}
\label{subsubsec:metriche_redflag}

Il modello di \textit{Red Flag Detection} affronta un task di classificazione binaria ad elevatissima criticità clinica. In questo contesto, il fallimento tecnico più grave è il \textbf{Falso Negativo}: non segnalare un paziente che manifesta reali indicatori di rischio suicidario. Le conseguenze di un mancato allarme sono, per definizione, irreversibili.

Per questa ragione, l'obiettivo primario è la \textbf{massimizzazione della Recall}, fissando un target minimo pari a:
\[
    \text{Recall} \geq 0{,}90
\]
ovvero il modello deve essere in grado di identificare correttamente almeno il 90\% dei casi effettivamente a rischio.

Il raggiungimento di una Recall elevata comporta inevitabilmente un numero maggiore di \textbf{Falsi Positivi} (pazienti segnalati come a rischio che in realtà non lo sono). Tuttavia, nel dominio clinico in esame, se da un lato un Falso Negativo può avere esiti catastrofici, dall'altro un eccesso di falsi allarmi produrrebbe un costo operativo insostenibile (troppe verifiche aggiuntive da parte dello psicologo) e il rischio di desensibilizzazione (\textit{alert fatigue}). Si definisce pertanto un limite inferiore rigoroso sulla Precision:
\[
    \text{Precision} \geq 0{,}70
\]
Tale soglia implica che, su 10 alert generati dal sistema, almeno 7 devono corrispondere a casi realmente critici, garantendo un eccellente bilanciamento tra la sicurezza del paziente e la sostenibilità operativa del carico di lavoro per gli psicologi.

Oltre al monitoraggio separato delle soglie di Recall e Precision, le performance globali del modello in fase di validazione verranno riassunte tramite l'$F_{\beta}$\textbf{-Score} (con $\beta > 1$). A differenza dell'$F_{1}$-Score, che bilancia equamente Precision e Recall, l'$F_{\beta}$-Score consente di pesare matematicamente l'importanza della Recall rispetto alla Precision di un fattore pari a $\beta^{2}$. Questa scelta riflette formalmente la criticità del dominio: pur esigendo un'alta affidabilità degli allarmi generati (Precision $\geq 0{,}70$), il costo clinico di un singolo caso non intercettato (Falso Negativo) rimane strutturalmente maggiore di quello associato a un'allerta superflua.

\subsubsection{Predizione PHQ-9 --- Regressione}
\label{subsubsec:metriche_phq9}

Il modello di \textit{Depression Prediction} è formulato come un task di regressione: stimare il punteggio numerico che il paziente otterrebbe compilando il questionario PHQ-9, su una scala da 0 a 27 \cite{kroenke2001phq9}. La metrica di valutazione primaria è il \textbf{Mean Absolute Error} (MAE), con un obiettivo massimo di:
\[
    \text{MAE} \leq 3 \text{ punti}
\]

La scelta della soglia di 3 punti è motivata dalla struttura clinica del PHQ-9, le cui fasce di severità hanno un'ampiezza di 5 punti (0--4 minima, 5--9 lieve, 10--14 moderata, 15--19 moderatamente grave, 20--27 grave). Un errore medio contenuto entro 3 punti accetta fisiologicamente il rumore e la soggettività insiti nei dati \textit{self-report}, ma preserva, nella stragrande maggioranza dei casi, la capacità del modello di collocare correttamente il paziente nella fascia di severità appropriata o in quella immediatamente adiacente.

Accanto al MAE, viene fissato un obiettivo sulla metrica \textbf{RMSE} (\textit{Root Mean Squared Error}):
\[
    \text{RMSE} \leq 4{,}5 \text{ punti}
\]
Poiché l'RMSE penalizza in modo quadratico gli scostamenti più ampi, il rispetto di questa soglia assicura che la \textit{baseline} predittiva generata dal modello rimanga stabile anche nei casi limite. Nel contesto del rilevamento delle compilazioni sospette, una predizione errata del modello non innesca un mancato intervento clinico, ma genera un ``falso allarme di incongruenza'', costringendo lo psicologo a revisionare manualmente un questionario che in realtà era veritiero e compilato in buona fede. La soglia di 4,5 punti --- volutamente inferiore all'ampiezza di una singola fascia di severità del PHQ-9 (5 punti) --- garantisce che il modello non produca mai errori isolati di magnitudo tale da invalidare ingiustamente le risposte oneste dei pazienti. Il monitoraggio congiunto di MAE e RMSE assicura pertanto un metro di paragone affidabile per l'intero processo di garanzia della qualità dei dati.

\subsubsection{Prevenzione Churn --- Ottimizzazione tramite Algoritmo Genetico}
\label{subsubsec:metriche_churn}

A differenza dei componenti precedenti, il sistema di prevenzione dell'abbandono non si configura come un task di classificazione o regressione, bensì come un \textbf{problema di ottimizzazione} risolto mediante un \textit{Algoritmo Genetico} (GA). L'obiettivo è evolvere strategie di \textit{nudging} --- combinazioni ottimali di tipo, frequenza e tempistica delle notifiche --- che massimizzino la probabilità di mantenere attivo il paziente sulla piattaforma. I criteri di successo tecnico si articolano pertanto su due assi complementari: la qualità della soluzione ottenuta e la performance algoritmica del processo evolutivo.

\paragraph{Qualità della Soluzione (Funzione di Fitness).}
La qualità di ciascuna strategia di \textit{nudging} candidata ($S$) viene quantificata attraverso una funzione di fitness formulata come combinazione lineare (somma pesata) dei singoli obiettivi:
\[
    \text{Fitness}(S) = w_1 \cdot \text{Retention}(S) - w_2 \cdot \text{Penalty\textsubscript{Freq}}(S) - w_3 \cdot \text{Penalty\textsubscript{Time}}(S)
\]
dove $\text{Retention}(S)$ stima il tasso di mantenimento atteso (coerente con l'obiettivo clinico $\geq 75\%$), $\text{Penalty\textsubscript{Freq}}$ modella una penalità esponenziale per l'invio eccessivo di notifiche (al fine di prevenire la \textit{notification fatigue}) e $\text{Penalty\textsubscript{Time}}$ penalizza l'invio in finestre temporali sub-ottimali o clinicamente sconsigliate. I pesi $w_1, w_2, w_3$ verranno calibrati empiricamente in fase di \textit{tuning} per garantire la priorità assoluta alla retention, imponendo al contempo soglie anti-spam ferree. L'algoritmo genetico ottimizzerà questo unico valore scalare per identificare la singola strategia migliore in modo diretto ed efficiente.

\paragraph{Performance Algoritmica.}
Affinché il GA possa essere impiegato come componente operativo nel \textit{backend} di SINTONIA, vengono fissati obiettivi quantitativi sul processo di ottimizzazione:
\begin{itemize}
    \item \textbf{Velocità di convergenza:} l'algoritmo deve convergere verso un \textit{plateau} di fitness soddisfacente entro un numero limitato di generazioni (indicativamente $G \leq 100$), dimostrando un'efficienza esplorativa idonea ad adattarsi ai cambiamenti comportamentali nel tempo.
    \item \textbf{Diversità genetica e prevenzione degli ottimi locali:} l'algoritmo deve mantenere un adeguato livello di diversità esplorativa durante le sue iterazioni, evitando una convergenza prematura verso sub-ottimi locali prima di aver scandagliato a fondo lo spazio delle soluzioni. Tale diversità viene monitorata quantificando costantemente la \textit{Distanza di Hamming} media tra i cromosomi della popolazione: l'obiettivo tecnico è che questo indicatore di varietà genetica non decada al di sotto di una soglia minima di guardia (es. il 15\% della distanza massima teorica) prima di aver raggiunto almeno la metà del budget generazionale prefissato.
\end{itemize}

% -- 1.3 ----------------------------------------------------------------------
\clearpage
\section{Specifiche PEAS degli Agenti Intelligenti}
\label{sec:peas_specifications}

L'architettura dei tre moduli di SINTON-IA è stata formalizzata secondo il framework dei sistemi intelligenti \textbf{PEAS} (\textit{Performance, Environment, Actuators, Sensors}). Questa schematizzazione non solo definisce i confini operativi degli agenti, ma ne chiarisce la responsabilità clinica e tecnica all'interno dell'ecosistema SINTONIA.

\begin{itemize}
    \item \textbf{Agente Diagnostico NLP (Red Flag Detection)} \\
          Un agente delegato alla sorveglianza proattiva dei contenuti testuali per l’identificazione di crisi imminenti.
          \begin{itemize}
              \item \textbf{Performance}: Massimizzazione della \textit{Recall} clinica ($\ge 0.90$) per minimizzare il rischio di falsi negativi, bilanciata da una \textit{Precision} $\ge 0.70$ per preservare l’attenzione operativa dello psicologo (\textit{alert fatigue}).
              \item \textbf{Environment}: Il database dei record clinici (diari) del paziente e il sistema di dashboard del personale sanitario preposto al triage.
              \item \textbf{Actuators}: Inserimento di un record strutturato (\texttt{alert\_clinico}) nel database del triage contenente il testo incriminato; generazione del flag di rischio per l'interfaccia UI.
              \item \textbf{Sensors}: Triggering asincrono \textit{event-driven} (fire-and-forget) attivato dal salvataggio di una nuova pagina del diario.
          \end{itemize}

          \vspace{0.3cm}
    \item \textbf{Agente Regressore (Depression Prediction)} \\
          Un agente analitico focalizzato sulla stima quantitativa del deterioramento dello stato mentale basato su marker emotivi.
          \begin{itemize}
              \item \textbf{Performance}: Contenimento del \textit{Mean Absolute Error} (MAE) $\le 3$ punti sulla scala PHQ-9; stabilità delle predizioni misurata tramite RMSE per limitare le divergenze ingiustificate tra auto-valutazione e predizione algoritmica.
              \item \textbf{Environment}: Il sistema informativo che traccia la storia clinica del paziente e l'applicativo di monitoraggio visionato dallo psicologo.
              \item \textbf{Actuators}: Scrittura dei trend PHQ-9 predetti su base dati; innesco logico di "Flag di Incongruenza" sui record tabellari sospetti.
              \item \textbf{Sensors}: Query periodiche per l'acquisizione delle serie temporali dei registri emotivi (\textit{Valence/Arousal}); parser dei log sulla frequenza di interazione quotidiana.
          \end{itemize}

          \vspace{0.3cm}
    \item \textbf{Agente Evolutivo (Churn Prevention)} \\
          Un agente ottimizzatore per personalizzare le strategie di engagement e prevenire l’abbandono terapeutico.
          \begin{itemize}
              \item \textbf{Performance}: Massimizzazione del tasso di \textit{User Retention} ($\ge 75\%$); minimizzazione delle penalità di \textit{Notification Fatigue} e rispetto dei vincoli circadiani degli utenti.
              \item \textbf{Environment}: L'utente finale inteso come sistema cognitivo-comportamentale, percepibile e modificabile unicamente attraverso l'infrastruttura di notifica mobile associata.
              \item \textbf{Actuators}: Modifica persistente delle policy di "Nudging" su database; schedulazione strutturata di payload contenenti template comunicativi, frequenze d'invio e finestre temporali.
              \item \textbf{Sensors}: Recettori di telemetria per i feedback impliciti; funzioni di calcolo aggregato sui log di adesione dell'utente.
          \end{itemize}
\end{itemize}

\subsection{Analisi del Contesto Operativo}
\label{subsec:task_environment}

Per completare la definizione degli agenti, è fondamentale mappare le proprietà dell’ambiente operativo in cui SINTON-IA agisce. L'integrazione di questi parametri permette di comprendere la complessità computazionale e l’incertezza con cui gli agenti devono confrontarsi.

\begin{table}[H]
    \centering
    \renewcommand{\arraystretch}{1.5}
    \small
    \resizebox{\textwidth}{!}{
    \begin{tabular}{|l|c|c|c|}
        \hline
        \textbf{Proprietà}     & \textbf{Agente NLP}      & \textbf{Agente Regressore} & \textbf{Agente Evolutivo} \\ \hline
        \textbf{Osservabilità} & Parzialmente Osservabile & Parzialmente Osservabile   & Parzialmente Osservabile  \\ \hline
        \textbf{Determinismo}  & Stocastico               & Stocastico                 & Stocastico                \\ \hline
        \textbf{Episodicità}   & Episodico                & Sequenziale                & Sequenziale               \\ \hline
        \textbf{Dinamismo}     & Statico                  & Statico                    & Dinamico                  \\ \hline
        \textbf{Discretezza}   & Discreto                 & Continuo                   & Discreto                  \\ \hline
        \textbf{Agenti}        & Singolo                  & Singolo                    & Singolo                   \\ \hline
    \end{tabular}
    }
    \caption{Proprietà dell'ambiente per i tre moduli di SINTON-IA.}
    \label{tab:task_environment}
\end{table}

L'ambiente è considerato \textbf{Parzialmente Osservabile} poiché l'agente non ha accesso allo stato mentale interno dell'utente, ma solo a manifestazioni digitali indirette (testo, log). È \textbf{Stocastico} a causa della natura imprevedibile del comportamento umano.

Il paradigma di esecuzione varia tra i moduli: l'Agente NLP opera in un ambiente \textbf{Episodico}, in quanto valuta il rischio leggendo esclusivamente l'ultimo testo inserito; al contrario, il Regressore e il GA sono \textbf{Sequenziali}, poiché analizzano lo storico dei dati del paziente (come l'andamento dell'umore o i giorni di inattività) per prendere decisioni.

Per quanto riguarda il \textbf{Dinamismo}, l'Agente NLP e il Regressore sono considerati \textbf{Statici}: lavorano su dati fermi nel database, che non cambiano nei pochi secondi in cui l'algoritmo effettua il calcolo. L'Agente Evolutivo affronta invece un ambiente \textbf{Dinamico}, dovendo adattare le notifiche a un utente reale le cui abitudini e livelli di attenzione variano continuamente durante la giornata.

Infine, per quanto riguarda la \textbf{Discretezza}: l'Agente NLP e l'Evolutivo lavorano in ambienti \textbf{Discreti}, poiché le loro decisioni ricadono su opzioni separate e ben distinte. L'ambiente del Regressore è invece \textbf{Continuo}, in quanto calcola l'andamento dell'umore tracciando variazioni fluide senza ricorrere a divisioni rigide.

\vspace{0.2cm}
Sebbene la Tabella~\ref{tab:task_environment} classifichi gli agenti come operanti in scenari ad \textbf{Agente Singolo} rispetto ai propri domini decisionali, si evidenzia una silente interazione pseudo-competitiva o collaborativa a livello sistemico: l'attività dell'Agente Evolutivo (modificando il grado di ingaggio dell'utente) altera la densità del flusso di dati che costituisce il retroterra sensorio dell'Agente Diagnostico e del Regressore, creando un ecosistema integrato.

\vspace{0.5cm}

\begin{figure}[H]
    \centering
    \includegraphics[width=\textwidth]{architettura-sintonia}
    \caption{Architettura Globale SINTON-IA: mappatura del flusso dati bidirezionale tra i moduli di Intelligenza Artificiale PEAS e il core applicativo React/NestJS del sistema padre.}
    \label{fig:architettura_globale}
\end{figure}

% -- 1.4 ----------------------------------------------------------------------
\section{Valutazione dei Rischi}
\label{sec:valutazione_rischi}

Il processo di integrazione dell'Intelligenza Artificiale all'interno dell'ecosistema SINTONIA comporta una serie di criticità intrinseche relative al trattamento dei dati sanitari, all'interazione tra modelli matematici e giudizio umano e alle variazioni operative nel tempo. Di seguito si analizzano i principali fattori di rischio associati ai modelli predittivi, suddivisi per ambito organizzativo e tecnologico, al fine di garantire l'efficacia a lungo termine e la sicurezza del sistema.

\subsection*{Rischi Tecnici}

Di seguito si analizzano le principali sfide di natura tecnica legate all'addestramento, all'affidabilità dei dati in ingresso e al ciclo di vita dei modelli di Machine Learning.

\begin{itemize}
    \item \textbf{Domain shift:} I modelli previsionali verranno addestrati sfruttando dataset provenienti da fonti pubbliche, le quali presentano inevitabilmente variazioni linguistiche, di distribuzione demografica e di contesto socioculturale rispetto alla reale popolazione dei pazienti in carico all'ASL Campania. Questa discrepanza rappresenta il rischio tecnico più critico dell'intero progetto (come evidenziato in merito alle risorse disponibili e ai loro limiti), potendo tradursi in una degradazione inattesa delle performance del modello in fase di produzione, nonostante valutazioni soddisfacenti durante il ciclo di addestramento e validazione.
    \item \textbf{Qualità dei dati in input (GIGO):} I pazienti in attesa potrebbero ingaggiare l'applicazione mobile in maniera intermittente o superficiale, arrivando a inserire stati d'animo casuali testuali brevi e ambigui o, nell'ipotesi più grave, a smettere progressivamente di compilare il diario. Tali comportamenti producono input inaffidabili che si ripercuotono a cascata sulle predizioni dei tre modelli, compromettendo irrimediabilmente la qualità del \textit{triage} dinamico (\textit{Garbage In, Garbage Out}).
    \item \textbf{Drift del modello nel tempo:} I modelli di Intelligenza Artificiale tendono a subire un calo fisiologico delle performance nel corso del tempo a causa dei mutamenti nel comportamento, nel linguaggio e nelle abitudini dell'utenza. Tale deragliamento statistico richiederà inevitabilmente la necessità di sessioni di re-training periodiche, introducendo un fattore di costo strutturale e di aggiornamento infrastrutturale già computato globalmente nelle stime OPEX dei requisiti di progetto.
\end{itemize}

\subsection*{Rischi Organizzativi}

L'efficacia del sistema non dipende unicamente dalla precisione algoritmica, ma anche dalla corretta interazione tra le previsioni dell'Intelligenza Artificiale e l'operatività quotidiana degli utenti finali (es. personale sanitario e pazienti in attesa).

\begin{itemize}
    \item \textbf{Alert fatigue:} Qualora il modello di Red Flag Detection per il rischio suicidario presenti, in scenario reale, un tasso di Falsi Positivi eccessivamente sproporzionato, l'elevato volume di allerte inattese finirebbe per desensibilizzare gli psicologi preposti alla gestione del triage. Tale fenomeno (\textit{Alert fatigue}) indurrebbe il personale sanitario a ignorare sistematicamente, o a declassare mentalmente, le notifiche di emergenza, vanificando la finalità stessa del modulo predittivo e il conseguimento del KPI di sostenibilità operativa stabilito nella Sezione \ref{sec:criteri_successo}.
    \item \textbf{Resistenza clinica:} Il personale psichiatrico e psicologico potrebbe sviluppare una diffidenza intrinseca verso le previsioni fornite dagli algoritmi (in particolar modo per il calcolo del trend depressivo PHQ-9 e del rischio suicidario), percependole non come un ausilio al filtro organizzativo ma come un'interferenza meccanica nel proprio giudizio diagnostico-clinico. Questa potenziale ostilità renderebbe estremamente complessa l'adozione e l'integrazione fluida degli applicativi nel panorama lavorativo reale delle ASL coinvolte.
    \item \textbf{Abbandono dell'app da parte dei pazienti:} Il monitoraggio e la valutazione dei vari rischi sono subalterni all'utilizzo stabile e continuativo della piattaforma. Qualora il tasso di abbandono (\textit{churn}) dei pazienti prima del primo incontro terapeutico raggiungesse livelli elevati, verrebbero a mancare l'intero insieme dei segnali clinici (diari e questionari) alla base dell'alimentazione dei tre modelli e dell'aggiornamento automatico della lista di priorità dei pazienti in attesa.
\end{itemize}




\subsection*{Piani di Contingenza}

\vspace{0.3cm}
\begin{table}[H]
    \centering
    \renewcommand{\arraystretch}{1.5}
    \small
    \begin{tabular}{|p{0.20\textwidth}|p{0.15\textwidth}|p{0.55\textwidth}|}
        \hline
        \textbf{Rischio}           & \textbf{Categoria} & \textbf{Piano di Contingenza}                                                                                                                                                                                                                                                                                                                                                                                                              \\
        \hline
        Domain shift               & Tecnico            & Validazione dei modelli su campioni rappresentativi del contesto italiano prima del deployment; previsione di una fase di fine-tuning su dati reali non appena disponibili.                                                                                                                                                                                                                                                                \\
        \hline
        GIGO                       & Tecnico            & Introduzione di meccanismi automatici di rilevamento delle anomalie nei pattern di compilazione (es. sequenze di valori identici o tempi di compilazione anomali).                                                                                                                                                                                                                                                                         \\
        \hline
        Drift del modello          & Tecnico            & Pianificazione di cicli periodici di monitoraggio delle performance e re-training, coerentemente con quanto previsto nelle voci OPEX.                                                                                                                                                                                                                                                                                                      \\
        \hline
        Alert fatigue              & Organizzativo      & Definizione di una soglia minima di confidenza per la generazione degli alert, calibrata sul KPI di sostenibilità operativa descritto nella relativa sezione.                                                                                                                                                                                                                                                                              \\
        \hline
        Resistenza clinica         & Organizzativo      & Progettazione di interfacce di output semplici e non invasive, orientate al supporto e non alla sostituzione del giudizio clinico.                                                                                                                                                                                                                                                                                                         \\
        \hline
        Abbandono dell'app         & Organizzativo      & Il Churn Prevention costituisce in sé stesso la misura preventiva primaria contro questo rischio.                                                                                                                                                                                                                                                                                                                                          \\
        \hline
    \end{tabular}
\end{table}
% -- 1.7 ----------------------------------------------------------------------

\section{Analisi Costi-Benefici e Valutazione delle Risorse}
\label{sec:costi_benefici}

Prima di avviare la fase operativa del progetto, è necessario condurre un'analisi strutturata delle risorse disponibili e dei costi associati allo sviluppo e al mantenimento dei modelli di Intelligenza Artificiale. Questa sezione mappa le risorse umane, dati e tecnologiche necessarie, quantifica i costi di sviluppo e gestione, e li confronta con i benefici attesi per l'ASL Campania, dimostrando la sostenibilità complessiva dell'iniziativa.


\clearpage
\subsection{Valutazione e Inventario delle Risorse Necessarie}
\label{subsec:inventario_risorse}

\subsubsection{Risorse Umane e Competenze}
\label{subsubsec:risorse_umane}

Il team di sviluppo di SINTON-IA è composto da tre membri --- A.W. De Fusco, A. Del Sorbo e G. Policola --- i quali si suddivideranno equamente le responsabilità secondo il modello organizzativo \textit{Team Data Science Process} (TDSP). Ciascun membro del team ricoprirà trasversalmente i ruoli previsti dal framework:

\begin{itemize}
    \item \textbf{Data Scientist:} progettazione, addestramento e validazione dei modelli predittivi;
    \item \textbf{Data Engineer:} costruzione delle pipeline di acquisizione, pulizia e trasformazione dei dati;
    \item \textbf{Sviluppatore Applicativo:} integrazione dei modelli nell'ecosistema SINTONIA tramite API di servizio.
\end{itemize}

Un aspetto critico riguarda la \textbf{validazione clinica} dei risultati prodotti dai modelli. Per garantire che le predizioni siano clinicamente significative e affidabili, sarà necessario il coinvolgimento futuro di \textit{esperti di dominio} --- psicologi clinici affiliati alle ASL campane --- che dovranno supervisionare l'etichettatura dei dati testuali (in particolare per il modello di Red Flag Detection) e validare la coerenza delle stime prodotte. Tuttavia, nella fase attuale del progetto, questa collaborazione non è ancora operativa: la validazione clinica formale è prevista come fase successiva, una volta consolidati i modelli di base mediante dataset pubblici.

\subsubsection{Risorse Dati e Vincoli di Privacy}
\label{subsubsec:risorse_dati}

L'acquisizione dei dati rappresenta una delle sfide più significative del progetto. Nel contesto europeo, i dati relativi alla salute mentale dei pazienti sono soggetti a rigide normative di tutela della privacy (GDPR, Regolamento UE 2016/679), che ne rendono di fatto impossibile l'accesso libero e gratuito per finalità di ricerca accademica non autorizzata.

Questa situazione configura un \textbf{dilemma strategico} che il progetto deve affrontare in modo consapevole:

\begin{itemize}
    \item \textbf{Dataset pubblici e liberamente accessibili} (es. piattaforme come Kaggle, repository accademici aperti): costituiscono la risorsa primaria per la fase di validazione iniziale dei modelli. Tali dataset sono tipicamente anonimizzati, provenienti da contesti extra-europei (prevalentemente anglofoni) e presentano inevitabili limiti di rappresentatività rispetto alla popolazione target campana. Le differenze linguistiche, culturali e cliniche possono influenzare le performance dei modelli, richiedendo opportune strategie di adattamento (es. traduzione, \textit{transfer learning}, \textit{fine-tuning} su dati più specifici).

    \item \textbf{Dataset per la ricerca clinica autorizzata}: rappresentano la risorsa ideale per affinare e rendere più accurati i modelli in una fase successiva del progetto. L'accesso a tali dati richiede la stipulazione di accordi formali con enti sanitari, l'approvazione di comitati etici e il rispetto di protocolli di anonimizzazione certificati. Questo percorso, pur essendo più oneroso in termini di tempo e burocrazia, garantirebbe dati clinicamente rilevanti e rappresentativi del contesto operativo reale.
\end{itemize}

La strategia adottata dal progetto prevede pertanto un \textbf{approccio incrementale}: nella fase iniziale si utilizzeranno esclusivamente dataset pubblici per la costruzione e la validazione di base dei modelli, posticipando a una fase successiva la richiesta di autorizzazioni formali laddove si rendesse necessario migliorare l'accuratezza e la completezza delle predizioni su dati più vicini al dominio di riferimento.

Per ciascun modello, le tipologie di dati necessarie sono riconducibili a:
\begin{itemize}
    \item \textbf{Red Flag Detection:} dati testuali annotati con etichette di rischio suicidario, contenenti testi redatti in prima persona che esprimano stati emotivi, ideazioni o segnali di disagio;
    \item \textbf{Depression Prediction:} serie temporali di indicatori emotivi (valenza, intensità) associati a punteggi di questionari clinici validati (PHQ-9 o equivalenti).
\end{itemize}

Per quanto riguarda il modello di \textbf{Churn Prevention}, non è previsto l'utilizzo di dataset esterni: il sistema si baserà esclusivamente sui dati interni generati dall'interazione degli utenti con la piattaforma SINTONIA (log di utilizzo, frequenza di accesso, pattern di compilazione), sfruttando l'approccio genetico-evolutivo per ottimizzare le strategie di intervento senza necessità di dati di addestramento etichettati.


\clearpage
\subsection{Analisi dei Costi (CAPEX e OPEX)}
\label{subsec:analisi_costi}

Nella presente fase, il progetto SINTON-IA si colloca interamente all'interno di un contesto accademico universitario: le risorse umane sono costituite da studenti, l'infrastruttura di calcolo è di proprietà del team e i dataset utilizzati sono pubblici e gratuiti. Di conseguenza, \textbf{i costi effettivi sostenuti per la prototipazione sono pari a zero}. Tuttavia, al fine di valutare la sostenibilità economica di un'eventuale adozione del sistema da parte dell'ASL Campania, è opportuno stimare i costi che si presenterebbero in uno scenario produttivo reale, suddividendoli nelle due categorie standard di analisi finanziaria: CAPEX (spese in conto capitale) e OPEX (spese operative ricorrenti).

\subsubsection{Costi di Sviluppo (CAPEX)}
\label{subsubsec:capex}

I costi CAPEX rappresentano gli investimenti iniziali necessari per portare il sistema dalla fase di prototipazione alla messa in produzione. In uno scenario di adozione reale da parte dell'ASL, le principali voci di costo sarebbero le seguenti:

\begin{itemize}
    \item \textbf{Risorse umane specializzate:} il coinvolgimento di figure professionali quali Data Scientist, Data Engineer e sviluppatori applicativi, ciascuno con competenze specifiche nel dominio della salute mentale e dell'Intelligenza Artificiale. A queste si aggiungerebbero le ore di consulenza di psicologi clinici per la supervisione e la validazione dei modelli.

    \item \textbf{Acquisizione e preparazione dei dati:} l'effort legato alla raccolta di dataset clinici autorizzati, alla loro anonimizzazione conforme al GDPR e alla fase di annotazione manuale (in particolare per i testi del diario necessari al modello di Red Flag Detection). Questa voce include anche eventuali costi legali per la stipulazione degli accordi con gli enti sanitari.

    \item \textbf{Infrastruttura di calcolo per l'addestramento:} il noleggio di risorse GPU necessarie per addestrare i modelli. A titolo indicativo, il noleggio di una GPU cloud di fascia alta ha un costo medio compreso tra 2 e 4 \texteuro/ora, rendendo l'addestramento di un singolo modello un investimento contenuto nell'ordine di poche centinaia di euro.

    \item \textbf{Integrazione nell'ecosistema SINTONIA:} lo sviluppo delle API di servizio, l'adattamento dell'interfaccia dell'applicazione per visualizzare gli output dei modelli (alert, grafici di andamento, notifiche strategiche) e i test di integrazione end-to-end.
\end{itemize}

\subsubsection{Costi Operativi e di Mantenimento (OPEX)}
\label{subsubsec:opex}

I costi OPEX rappresentano le spese ricorrenti necessarie per mantenere il sistema operativo una volta in produzione. A differenza dei costi CAPEX, che vengono sostenuti una sola volta, gli OPEX si ripetono con cadenza periodica:

\begin{itemize}
    \item \textbf{Costi di inferenza:} ogni volta che un paziente scrive una pagina di diario, il modello di Red Flag Detection deve analizzare il testo in tempo reale. Analogamente, le predizioni PHQ-9 e le valutazioni di churn richiedono elaborazioni periodiche. Questi costi sono proporzionali al numero di utenti attivi e alla frequenza di utilizzo della piattaforma. Utilizzando servizi cloud, il costo unitario per singola inferenza è nell'ordine di frazioni di centesimo, ma può accumularsi con volumi elevati di utenti.

          % Tradurre gli obiettivi di business in obiettivi tecnici di data mining:
          % - Classificazione del rischio suicidario tramite NLP
          % - Prevenzione del churn terapeutico
          % - Rilevamento automatico della depressione
          % - Integrazione dei modelli tramite API

    \item \textbf{Ri-addestramento periodico:} per contrastare il data drift e incorporare nuovi dati clinici, i modelli dovranno essere periodicamente ri-addestrati. Ogni ciclo di ri-addestramento comporta costi di calcolo GPU e di validazione clinica da parte degli esperti di dominio.
\end{itemize}

\clearpage
\subsection{Valutazione dei Benefici e Ritorno sull'Investimento}
\label{subsec:benefici_roi}

La sostenibilità di un progetto come SINTON-IA non può essere valutata esclusivamente in termini economici: il valore primario del sistema risiede nella sua capacità di \textbf{salvare vite umane} e di trasformare radicalmente la qualità del monitoraggio psicologico durante una delle fasi più critiche e vulnerabili del percorso di cura. I benefici attesi si articolano su tre livelli complementari.

\subsubsection{Benefici Clinici: Monitoraggio Continuo e Prevenzione del Rischio}
\label{subsubsec:benefici_clinici}

Il beneficio più significativo e non negoziabile dell'intero progetto è la capacità di \textbf{intercettare tempestivamente situazioni di rischio suicidario} che, nel sistema attuale, resterebbero invisibili per giorni o settimane.

Come evidenziato nella Sezione~\ref{sec:flusso_psicologo}, il sistema di alert attuale è intrinsecamente limitato dalla periodicità dei questionari: un paziente che esprime segnali di pericolo nel proprio diario nei giorni che intercorrono tra una somministrazione e l'altra non viene identificato. SINTON-IA colma questo vuoto critico attraverso il modello di \textit{Red Flag Detection}, che analizza in tempo reale e in modo completamente anonimo ogni pagina scritta nel diario, generando un alert clinico immediato al primo riscontro di pattern linguistici riconducibili a ideazione suicidaria.

L'impatto clinico di questa funzionalità è potenzialmente incommensurabile: \textbf{anche un solo caso intercettato in tempo può rappresentare la differenza tra la vita e la morte di un paziente}. In termini economici, la prevenzione di un singolo episodio di emergenza psichiatrica acuta --- che comporta ospedalizzazione d'urgenza, degenza in reparto psichiatrico, trattamenti farmacologici intensivi e follow-up prolungato --- genera un risparmio per il Sistema Sanitario Regionale stimabile in diverse migliaia di euro per episodio evitato. Tuttavia, ridurre il valore di questa funzionalità a un mero calcolo economico sarebbe profondamente riduttivo: il vero ritorno sull'investimento è la protezione della persona.

Analogamente, il modello di \textit{Depression Prediction} estende il monitoraggio ben oltre le finestre temporali dei questionari, offrendo allo psicologo una visione continua e granulare dell'andamento emotivo del paziente. La possibilità di rilevare un peggioramento del quadro depressivo \textit{prima} che si manifesti nella compilazione successiva del PHQ-9 consente interventi preventivi che possono evitare l'aggravamento della condizione clinica, riducendo il rischio di cronicizzazione e i costi sanitari ad essa associati.

\subsubsection{Benefici Operativi: Supporto alla Supervisione dello Psicologo}
\label{subsubsec:benefici_operativi}

Nella versione attuale della piattaforma, le possibilità operative dello psicologo sono circoscritte: il professionista può supervisionare i questionari compilati dai pazienti in attesa, accettare o gestire gli alert clinici generati dal sistema e, una volta effettuata la presa in carico, il percorso del paziente sulla piattaforma si conclude. All'interno di questo perimetro, SINTON-IA introduce due benefici operativi concreti.

In primo luogo, il modello di \textit{Depression Prediction} trasforma radicalmente la \textbf{fase di revisione dei questionari}. Attualmente lo psicologo è costretto a procedere con controlli a campione casuali, sperando di imbattersi in compilazioni palesemente incongruenti --- un processo lento, dispendioso e con basse probabilità di successo. Grazie alla predizione del punteggio PHQ-9 basata sugli stati d'animo registrati dal paziente, il sistema può confrontare automaticamente il valore atteso con quello effettivamente dichiarato nel questionario, \textbf{segnalando in modo mirato le compilazioni sospette} anziché costringere il professionista a cercarle alla cieca. Lo psicologo potrà così dedicare il proprio tempo esclusivamente alla verifica dei casi evidenziati, con un'efficacia incomparabilmente superiore.

In secondo luogo, le predizioni del modello offrono allo psicologo una \textbf{visione continua dell'andamento emotivo} di ciascun paziente. Attraverso i grafici generati dalla stima quotidiana del punteggio PHQ-9, il professionista può osservare l'evoluzione dell'umore nel tempo, individuare tendenze al peggioramento e disporre di un quadro clinico più ricco al momento della presa in carico. Questa informazione, altrimenti inaccessibile nei lunghi intervalli tra una somministrazione e l'altra, consente decisioni più informate e tempestive.

\subsubsection{Benefici di Servizio: Retention e Qualità dell'Assistenza Pubblica}
\label{subsubsec:benefici_servizio}

Il terzo livello di benefici riguarda la qualità complessiva del servizio offerto dal sistema sanitario pubblico campano. Uno dei problemi più insidiosi dell'attuale modello di attesa è l'\textbf{abbandono silenzioso}: pazienti che, scoraggiati dall'assenza di interazione e dalla percezione di abbandono istituzionale, smettono di utilizzare la piattaforma o rinunciano del tutto alle cure.

Il modello di \textit{Churn Prevention} affronta direttamente questo problema, identificando i pazienti a rischio di abbandono e attivando strategie di \textit{nudging} personalizzate --- notifiche mirate, promemoria, contenuti di supporto --- che trasformano il periodo di attesa da una fase passiva e frustrante in un percorso di \textbf{presa in carico attiva}.

I benefici di questa trasformazione sono molteplici:
\begin{itemize}
    \item la continuità nell'utilizzo della piattaforma garantisce la raccolta costante dei dati comportamentali, alimentando il triage dinamico e migliorando l'accuratezza di tutti i modelli predittivi;
    \item il paziente percepisce un'attenzione concreta e personalizzata da parte del servizio sanitario, riducendo il senso di abbandono e rafforzando la fiducia nell'istituzione;
    \item l'ASL migliora significativamente la propria immagine pubblica, dimostrando di investire in strumenti innovativi e di adottare un approccio proattivo alla salute mentale dei cittadini.
\end{itemize}

In sintesi, il ritorno sull'investimento di SINTON-IA non si misura esclusivamente in termini economici, ma nella capacità del sistema di \textbf{rendere il servizio di supporto psicologico pubblico più reattivo, più umano e, soprattutto, più sicuro} per i cittadini campani in attesa di presa in carico.
